\chapter{2009年陕西卷（理）}

\section{单选题}

\begin{problem}
设不等式 \(x^{2} - x \leqslant 0\) 的解集为 \(M\)，函数 \(f\left(x\right) = \ln\left(1 - |x|\right)\) 的定义域为 \(N\)，则 \(M \cap N\) 为（\quad）

\choices
    {\(\left[0, 1\right)\)}
    {\(\left(0,1\right)\)}
    {\(\left[0,1\right]\)}
    {\(\left(-1,0\right]\)}

\begin{answer}
A
\end{answer}
\end{problem}

\begin{problem}
已知 \(z\) 是纯虚数，\(\frac{z + 2}{1 - \i}\) 是实数，那么 \(z\) 等于（\quad）

\choices
    {\(2\i\)}
    {\(\i\)}
    {\(-\i\)}
    {\(-2\i\)}

\begin{answer}
D
\end{answer}
\end{problem}

\begin{problem}
函数 \( f\left(x\right) = \sqrt{2x - 4} \)（\(x \geqslant 4\)）的反函数为（\quad）

\choices
    {\(f^{-1}\left(x\right) = \frac{1}{2} x^2 + 2\)（\(x\geqslant 0\)）}
    {\(f^{-1}\left(x\right) = \frac{1}{2} x^2 + 2\)（\(x\geqslant 2\)）}
    {\(f^{-1}\left(x\right) = \frac{1}{2} x^2 + 4\)（\(x\geqslant 0\)）}
    {\(f^{-1}\left(x\right) = \frac{1}{2} x^2 + 4\)（\(x\geqslant 2\)）}

\begin{answer}
B
\end{answer}
\end{problem}

\begin{problem}
过原点且倾斜角为 \(60^{\circ}\) 的直线被圆 \(x^{2} + y^{2} - 4y = 0\) 所截得的弦长为（\quad）

\choices
    {\(\sqrt{3}\)}
    {\(2\)}
    {\( \sqrt{6} \)}
    {\(2\sqrt{3}\)}

\begin{answer}
D
\end{answer}
\end{problem}

\begin{problem}
若 \(3\sin \alpha + \cos \alpha = 0\)，则 \(\frac{1}{\cos^2\alpha + \sin 2\alpha}\) 的值为（\quad）

\choices
    {\(\frac{10}{3}\)}
    {\( \frac{5}{3} \)}
    {\( \frac{2}{3} \)}
    {\(-2\)}

\begin{answer}
A
\end{answer}
\end{problem}

\begin{problem}
若 \( \left(1-2x\right)^{2009}=a_{0}+a_{1}x+\cdots+a_{2009}x^{2009} \) （\( x\in \R \)），则 \( \frac{a_{1}}{2}+\frac{a_{2}}{2^{2}}+\cdots+\frac{a_{2009}}{2^{2009}} \)（\quad）

\choices
    {\(2\)}
    {\(0\)}
    {\(-1\)}
    {\(-2\)}

\begin{answer}
C
\end{answer}
\end{problem}

\begin{problem}
“ \( m > n > 0 \) ”是“方程 \( mx^{2} + ny^{2} = 1 \) 表示焦点在 \(y\) 轴上的椭圆”的（\quad）

\choices
    {充分而不必要条件}
    {必要而不充分条件}
    {充要条件}
    {既不充分也不必要条件}

\begin{answer}
C
\end{answer}
\end{problem}

\begin{problem}
在 \(\triangle ABC\) 中，\(M\) 是 \(BC\) 的中点，\(AM = 1\)，点 \(P\) 在 \(AM\) 上且满足 \(\overrightarrow{AP} = 2\overrightarrow{PM}\)，则 \(\overrightarrow{PA} \cdot \left(\overrightarrow{PB} + \overrightarrow{PC}\right)\) 等于（\quad）

\choices
    {\(-\frac{4}{9}\)}
    {\(-\frac{4}{3}\)}
    {\( \frac{4}{3} \)}
    {\( \frac{4}{9} \)}

\begin{answer}
A
\end{answer}
\end{problem}

\begin{problem}
从 \(0\)，\(1\)，\(2\)，\(3\)，\(4\)，\(5\) 这六个数中任取两个奇数和两个偶数，组成没有重复数字的四位数的个数为（\quad）

\choices
    {\(300\)}
    {\(216\)}
    {\(180\)}
    {\(162\)}

\begin{answer}
C
\end{answer}
\end{problem}

\begin{problem}
若正方体的棱长为 \(\sqrt{2}\)，则以该正方体各个面的中心为顶点的凸多面体的体积为（\quad）

\choices
    {\(\frac{\sqrt{2}}{6}\)}
    {\(\frac{\sqrt{2}}{3}\)}
    {\(\frac{\sqrt{3}}{3}\)}
    {\(\frac{2}{3}\)}

\begin{answer}
B
\end{answer}
\end{problem}

\begin{problem}
若 \(x\)，\(y\) 满足约束条件 \(\begin{cases}
x + y \geqslant 1,\\
x - y \geqslant -1,\\
2x - y \leqslant 2,
\end{cases}\) 目标函数 \(z = ax + 2y\) 仅在点 \(\left(1,0\right)\) 处取得最小值，则 \(a\) 的取值范围是（\quad）

\choices
    {\(\left(-1, 2\right)\)}
    {\(\left(-4, 2\right)\)}
    {\(\left(-4, 0\right]\)}
    {\(\left(-2, 4\right)\)}

\begin{answer}
B
\end{answer}
\end{problem}

\begin{problem}
定义在 \(\R\) 上的偶函数 \(f\left(x\right)\) 满足：对任意的 \(x_{1}, x_{2} \in \left(-\infty, 0\right]\) （\(x_{1} \neq x_{2}\)），有 \(\left(x_{2} - x_{1}\right)\left(f\left(x_{2}\right) - f\left(x_{1}\right)\right) > 0\). 则当 \(n \in \N^{*}\) 时，有（\quad）

\choices
    {\(f\left(-n\right) < f\left(n - 1\right) < f\left(n + 1\right)\)}
    {\(f\left(n - 1\right) < f\left(-n\right) < f\left(n + 1\right)\)}
    {\(f\left(n + 1\right) < f\left(-n\right) < f\left(n - 1\right)\)}
    {\(f\left(n + 1\right) < f\left(n - 1\right) < f\left(-n\right)\)}

\begin{answer}
C
\end{answer}
\end{problem}

\section{填空题}

\begin{problem}
设等差数列 \(\{a_{n}\}\) 的前 \(n\) 项和为 \(S_{n}\)，若 \(a_{6} = S_{3} = 12\)，则 \(\displaystyle\lim_{n\to \infty}\frac{S_n}{n^2} =\)\fillinblank{}

\begin{answer}
\(1\)
\end{answer}
\end{problem}

\begin{problem}
某班有 \(36\) 名同学参加数学，物理，化学课外探究小组，每名同学至多参加两个小组，已知参加数学，物理，化学小组的人数分别为 \(26\)，\(15\)，\(13\)，同时参加数学和物理小组的有 \(6\) 人，同时参加物理和化学小组的有 \(4\) 人，同时参加数学和化学小组的有\fillinblank{}人.

\begin{answer}
\(8\)
\end{answer}
\end{problem}

\begin{problem}
如图球 \(O\) 的半径为 \(2\)，圆 \(O_{1}\) 是一小圆，\(O_{1}O = \sqrt{2}\)，\(A\)，\(B\) 是圆 \(O_{1}\) 上两点. 若 \(A\)，\(B\) 两点间的球面距离为 \(\frac{2\pi}{3}\)，则 \(\angle AO_{1}B =\fillinblank{}\).

\bitmapfigure[max width=0.34\linewidth]{img/2009/shaanxi_science/q15_fig1.png}

\begin{answer}
\(\frac{\pi}{2}\)
\end{answer}
\end{problem}

\begin{problem}
设曲线 \( y = x^{n + 1} \)（\(n \in \N^*\)）在点 \(\left(1,1\right)\) 处的切线与 \(x\) 轴的交点的横坐标为 \( x_n \). 令 \( a_n = \lg x_n \)，则 \( a_1 + a_2 + \dots + a_{99} \) 的值为\fillinblank{}.

\begin{answer}
\(-2\)
\end{answer}
\end{problem}

\section{解答题}

\begin{problem}
已知函数 \(f\left(x\right) = A \sin\left(\omega x + \varphi\right)\)，\(x \in \R\) （其中 \(A > 0\)，\(\omega > 0\)，\(0 < \varphi < \frac{\pi}{2}\)）的图象与 \(x\) 轴的交点中，相邻两个交点之间的距离为 \( \frac{\pi}{2} \)，且图象上一个最低点为 \( M\left(\frac{2\pi}{3}, -2\right) \).

\begin{enumerate}
\item 求 \( f\left(x\right) \) 的解析式；
\item 当 \(x \in \left[\frac{\pi}{12}, \frac{\pi}{2}\right]\)，求 \(f\left(x\right)\) 的值域.
\end{enumerate}

\end{problem}

\begin{solution}
\begin{enumerate}
\item 函数图象与 \(x\) 轴的交点满足 \(\sin\left(\omega x+\varphi\right)=0\)，相邻两个交点对应的相位相差 \(\pi\)，所以相邻交点的横坐标之差满足 \(\omega\cdot\frac{\pi}{2}=\pi\)，从而 \(\omega=2\). 由最低点的纵坐标为 \(-2\)，得振幅 \(A=2\). 当 \(x=\frac{2\pi}{3}\) 时函数取得最小值，因此 \(2\cdot\frac{2\pi}{3}+\varphi=\frac{3\pi}{2}+2k\pi\)，其中 \(k\in\Z\). 于是 \(\varphi=\frac{\pi}{6}+2k\pi\)，结合 \(0<\varphi<\frac{\pi}{2}\) 得 \(\varphi=\frac{\pi}{6}\)，所以 \(f\left(x\right)=2\sin\left(2x+\frac{\pi}{6}\right)\).
\item 令 \(t=2x+\frac{\pi}{6}\)，则当 \(x\in\left[\frac{\pi}{12},\frac{\pi}{2}\right]\) 时，\(t\in\left[\frac{\pi}{3},\frac{7\pi}{6}\right]\). 在此区间内，\(\sin t\) 在 \(t=\frac{\pi}{2}\) 处取得最大值 \(1\)，在右端点 \(t=\frac{7\pi}{6}\) 处取得最小值 \(-\frac{1}{2}\). 因而 \(f\left(x\right)=2\sin t\) 的值域为 \(\left[-1,2\right]\).
\end{enumerate}
\end{solution}

\begin{problem}
如图，在直三棱柱 \( ABC - A_{1}B_{1}C_{1} \) 中，\( AB = 1 \)，\( AC = AA_{1} = \sqrt{3} \)，\( \angle ABC = 60^{\circ} \).

\begin{enumerate}
\item 证明： \( AB \perp A_{1}C \)；
\item 求二面角 \(A - A_{1}C - B\) 的大小.
\end{enumerate}

\bitmapfigure[max width=0.42\linewidth]{img/2009/shaanxi_science/q18_fig1.png}
\end{problem}

\begin{solution}
\begin{enumerate}
\item 在 \(\triangle ABC\) 中，由余弦定理得 \(AC^{2}=AB^{2}+BC^{2}-2AB\cdot BC\cos60^{\circ}\)，即 \(3=1+BC^{2}-BC\). 解得 \(BC=2\)（舍去负根），于是 \(AC^{2}+AB^{2}=3+1=4=BC^{2}\)，所以 \(\angle BAC=90^{\circ}\)，即 \(AB\perp AC\). 又因为三棱柱是直三棱柱，\(AA_{1}\perp\) 平面 \(ABC\)，从而 \(AB\perp AA_{1}\). 直线 \(AC\) 与 \(AA_{1}\) 是平面 \(AA_{1}C\) 内相交的两条直线，且 \(AB\) 同时垂直于它们，所以 \(AB\perp\) 平面 \(AA_{1}C\)，进而 \(AB\perp A_{1}C\).
\item 建立如图所示的直角坐标系，取 \(A=\left(0,0,0\right)\)，\(C=\left(\sqrt{3},0,0\right)\)，\(B=\left(0,1,0\right)\)，\(A_{1}=\left(0,0,\sqrt{3}\right)\). 则 \(\overrightarrow{A_{1}C}=\left(\sqrt{3},0,-\sqrt{3}\right)\). 设 \(H\) 为 \(A_{1}C\) 的中点，则 \(H=\left(\frac{\sqrt{3}}{2},0,\frac{\sqrt{3}}{2}\right)\)，并且 \(\overrightarrow{HA}=\left(-\frac{\sqrt{3}}{2},0,-\frac{\sqrt{3}}{2}\right)\)，\(\overrightarrow{HB}=\left(-\frac{\sqrt{3}}{2},1,-\frac{\sqrt{3}}{2}\right)\). 直接计算得 \(\overrightarrow{HA}\cdot\overrightarrow{A_{1}C}=0\)，\(\overrightarrow{HB}\cdot\overrightarrow{A_{1}C}=0\)，所以 \(HA\) 和 \(HB\) 分别是两个半平面 \(A-A_{1}C\)、\(B-A_{1}C\) 内垂直于棱 \(A_{1}C\) 的射线. 设二面角 \(A-A_{1}C-B\) 为 \(\theta\)，则
\[
\cos\theta=\frac{\overrightarrow{HA}\cdot\overrightarrow{HB}}{|\overrightarrow{HA}|\,|\overrightarrow{HB}|}=\frac{\frac{3}{2}}{\sqrt{\frac{3}{2}}\sqrt{\frac{5}{2}}}=\frac{3}{\sqrt{15}}=\frac{\sqrt{15}}{5}.
\]
因此二面角 \(A-A_{1}C-B\) 的大小为 \(\arccos\frac{\sqrt{15}}{5}\).
\end{enumerate}
\end{solution}

\begin{problem}
某食品企业一个月内被消费者投诉的次数用 \(\xi\) 表示，据统计，随机变量 \(\xi\) 的概率分布如下：

\begin{center}
\begin{tabular}{|c|c|c|c|c|}
\hline
\(\xi\) & \(0\) & \(1\) & \(2\) & \(3\) \\
\hline
\(p\) & \(0.1\) & \(0.3\) & \(2a\) & \(a\) \\
\hline
\end{tabular}
\end{center}

\begin{enumerate}
\item 求 \(a\) 的值和 \(\xi\) 的数学期望；
\item 假设一月份与二月份被消费者投诉的次数互不影响，求该企业在这两个月内共被消费者投诉 2 次的概率.
\end{enumerate}
\end{problem}

\begin{solution}
\begin{enumerate}
\item 概率之和为 \(1\)，所以 \(0.1+0.3+2a+a=1\)，解得 \(a=0.2\). 因而 \(\xi\) 的数学期望为 \(E\left(\xi\right)=0\cdot0.1+1\cdot0.3+2\cdot\left(2a\right)+3\cdot a=0.3+0.8+0.6=1.7\).
\item 设一月份和二月份的投诉次数分别为 \(\xi_{1}\) 和 \(\xi_{2}\). 由于两个月互不影响，\(\xi_{1}\) 与 \(\xi_{2}\) 相互独立. 两个月共被投诉 \(2\) 次包括 \(\left(\xi_{1},\xi_{2}\right)=\left(0,2\right)\)、\(\left(1,1\right)\)、\(\left(2,0\right)\) 三种互斥情形，所以所求概率为 \(0.1\cdot0.4+0.3\cdot0.3+0.4\cdot0.1=0.17\).
\end{enumerate}
\end{solution}

\begin{problem}
已知函数 \(f\left(x\right) = \ln\left(ax + 1\right) + \frac{1 - x}{1 + x}\)，\(x \geqslant 0\)，其中 \( a > 0 \).

\begin{enumerate}
\item 若 \( f\left(x\right) \) 在 \( x = 1 \) 处取得极值，求 \( a \) 的值；
\item 求 \( f\left(x\right) \) 的单调区间；
\item 若 \( f\left(x\right) \) 的最小值为 1，求 \(a\) 的取值范围.
\end{enumerate}
\end{problem}

\begin{solution}
\begin{enumerate}
\item 函数的导数为 \(f'\left(x\right)=\frac{a}{ax+1}-\frac{2}{\left(1+x\right)^{2}}\). 因为 \(x=1\) 是定义域的内点，若函数在此处取得极值，则必有 \(f'\left(1\right)=0\)，即 \(\frac{a}{a+1}-\frac{1}{2}=0\)，解得 \(a=1\). 当 \(a=1\) 时，\(f'\left(x\right)=\frac{x-1}{\left(1+x\right)^{2}}\)，所以函数在 \(x=1\) 的左侧递减、右侧递增，确实在 \(x=1\) 处取得极小值.
\item 因为 \(a>0\) 且 \(x\geqslant0\)，有
\[
f'\left(x\right)=\frac{a\left(1+x\right)^{2}-2\left(1+ax\right)}{\left(1+ax\right)\left(1+x\right)^{2}}=\frac{ax^{2}+a-2}{\left(1+ax\right)\left(1+x\right)^{2}}.
\]
分母恒为正，故只需讨论 \(ax^{2}+a-2\) 的符号. 当 \(0<a<2\) 时，令 \(x_{0}=\sqrt{\frac{2-a}{a}}\)，则 \(f'\left(x\right)<0\)（\(0\leqslant x<x_{0}\)），\(f'\left(x\right)>0\)（\(x>x_{0}\)），所以函数在 \(\left[0,x_{0}\right]\) 上单调递减，在 \(\left[x_{0},+\infty\right)\) 上单调递增. 当 \(a\geqslant2\) 时，\(ax^{2}+a-2\geqslant0\)，所以函数在 \([0,+\infty)\) 上单调递增.
\item 当 \(a\geqslant2\) 时，由上问知函数在 \([0,+\infty)\) 上单调递增，且 \(f\left(0\right)=1\)，故函数的最小值为 \(1\). 下面考察 \(0<a<2\) 的情形. 此时函数在 \(x_{0}=\sqrt{\frac{2-a}{a}}\) 处取得最小值，令 \(t=x_{0}>0\)，则 \(a=\frac{2}{1+t^{2}}\). 于是
\[
f\left(t\right)-1=\ln\left(1+\frac{2t}{1+t^{2}}\right)-\frac{2t}{1+t}=\ln\frac{\left(1+t\right)^{2}}{1+t^{2}}-\frac{2t}{1+t}.
\]
记 \(h\left(t\right)=\ln\frac{\left(1+t\right)^{2}}{1+t^{2}}-\frac{2t}{1+t}\)，则 \(h\left(0\right)=0\)，且
\[
h'\left(t\right)=\frac{2}{1+t}-\frac{2t}{1+t^{2}}-\frac{2}{\left(1+t\right)^{2}}=-\frac{4t^{2}}{\left(1+t\right)^{2}\left(1+t^{2}\right)}<0 \quad \left(t>0\right).
\]
所以 \(h\left(t\right)<0\)，即此时函数的最小值小于 \(1\). 因此函数的最小值为 \(1\) 时，\(a\) 的取值范围为 \([2,+\infty)\).
\end{enumerate}
\end{solution}

\begin{problem}
已知双曲线 \(C\) 的方程为 \( \frac{y^{2}}{a^{2}} - \frac{x^{2}}{b^{2}} = 1 \)（\(a > 0\)，\(b > 0\)），离心率 \( e = \frac{\sqrt{5}}{2} \)，顶点到渐近线的距离为 \( \frac{2\sqrt{5}}{5} \).

\begin{enumerate}
\item 求双曲线 \(C\) 的方程；
\item 如图，\(P\) 是双曲线 \(C\) 上一点，\(A\)，\(B\) 两点在双曲线 \(C\) 的两条渐近线上，且分别位于第一、二象限. 若 \(\overrightarrow{AP} = \lambda \overrightarrow{PB}\)，\(\lambda \in \left[\frac{1}{3}, 2\right]\)，求 \(\triangle AOB\) 面积的取值范围.
\end{enumerate}

\bitmapfigure[max width=0.42\linewidth]{img/2009/shaanxi_science/q21_fig1.png}
\end{problem}

\begin{solution}
\begin{enumerate}
\item 设双曲线的焦半距为 \(c\)，则 \(c^{2}=a^{2}+b^{2}\)，且由离心率得 \(\frac{c}{a}=\frac{\sqrt{5}}{2}\)，所以 \(\frac{b^{2}}{a^{2}}=\frac{c^{2}-a^{2}}{a^{2}}=\frac{1}{4}\)，从而 \(b=\frac{a}{2}\). 双曲线的渐近线为 \(y=\pm\frac{a}{b}x=\pm2x\). 顶点 \(\left(0,a\right)\) 到直线 \(y=2x\) 的距离为 \(\frac{a}{\sqrt{5}}\)，由题意 \(\frac{a}{\sqrt{5}}=\frac{2\sqrt{5}}{5}=\frac{2}{\sqrt{5}}\)，得 \(a=2\)，\(b=1\). 因此双曲线 \(C\) 的方程为 \(\frac{y^{2}}{4}-x^{2}=1\).
\item 设 \(P=\left(x,y\right)\). 由图中位置，\(P\) 在双曲线的上支，故 \(y>2|x|\). 设第一象限内渐近线上的点为 \(A=\left(u,2u\right)\)，第二象限内另一条渐近线上的点为 \(B=\left(-v,2v\right)\)，其中 \(u>0\)，\(v>0\). 由 \(\overrightarrow{AP}=\lambda\overrightarrow{PB}\) 得 \(\left(1+\lambda\right)P=A+\lambda B\)，于是
\[
\left(1+\lambda\right)x=u-\lambda v,
\qquad \left(1+\lambda\right)y=2u+2\lambda v.
\]
解得 \(u=\frac{1+\lambda}{4}\left(2x+y\right)\)，\(v=\frac{1+\lambda}{4\lambda}\left(y-2x\right)\). 因而
\[
S_{\triangle AOB}=\frac{1}{2}\left|\begin{vmatrix}u&2u\\-v&2v\end{vmatrix}\right|=2uv=\frac{\left(1+\lambda\right)^{2}}{8\lambda}\left(y^{2}-4x^{2}\right)=\frac{\left(1+\lambda\right)^{2}}{2\lambda},
\]
其中最后一步使用了 \(P\in C\)，即 \(y^{2}-4x^{2}=4\). 记 \(S\left(\lambda\right)=\frac{\left(1+\lambda\right)^{2}}{2\lambda}\)，则 \(S'\left(\lambda\right)=\frac{\lambda^{2}-1}{2\lambda^{2}}\)，所以 \(S\left(\lambda\right)\) 在 \(\left[\frac{1}{3},1\right]\) 上递减，在 \([1,2]\) 上递增. 又 \(S\left(1\right)=2\)，\(S\left(\frac{1}{3}\right)=\frac{8}{3}\)，\(S\left(2\right)=\frac{9}{4}\)，故 \(\triangle AOB\) 面积的取值范围为 \(\left[2,\frac{8}{3}\right]\).
\end{enumerate}
\end{solution}

\begin{problem}
已知数列 \(\{x_{n}\}\) 满足 \(x_{1} = \frac{1}{2}\)，\(x_{n + 1} = \frac{1}{1 + x_n}\)，\(n \in \N^*\).

\begin{enumerate}
\item 猜想数列 \( \{x_{2n}\} \) 的单调性，并证明你的结论；
\item 证明： \( \left|x_{n+1}-x_{n}\right|\leqslant\frac{1}{6}\left(\frac{2}{5}\right)^{n-1} \).
\end{enumerate}
\end{problem}

\begin{solution}
\begin{enumerate}
\item 记 \(g\left(x\right)=\frac{1}{1+\frac{1}{1+x}}=\frac{1+x}{2+x}\)，则 \(x_{n+2}=g\left(x_n\right)\)，且 \(g\left(x\right)-x=\frac{1-x-x^{2}}{2+x}\).
令 \(\alpha=\frac{\sqrt{5}-1}{2}\)，则 \(1-\alpha-\alpha^{2}=0\)，所以当 \(x>0\) 时，\(g\left(x\right)>x\)（\(0<x<\alpha\)），\(g\left(x\right)<x\)（\(x>\alpha\)）. 又 \(x_{2}=\frac{2}{3}>\alpha\)，并且 \(g'\left(x\right)=\frac{1}{\left(2+x\right)^{2}}>0\). 因此只要 \(x_{2n}>\alpha\)，就有 \(x_{2n+2}=g\left(x_{2n}\right)<x_{2n}\)，同时 \(x_{2n+2}>g\left(\alpha\right)=\alpha\). 归纳可得 \(x_{2n}>\alpha\) 且 \(x_{2n+2}<x_{2n}\)，所以数列 \(\{x_{2n}\}\) 单调递减.
\item 由 \(x_{1}=\frac{1}{2}\) 以及递推关系可归纳得到 \(\frac{1}{2}\leqslant x_n\leqslant\frac{2}{3}\)：若 \(\frac{1}{2}\leqslant x_n\leqslant\frac{2}{3}\)，则 \(\frac{3}{5}\leqslant x_{n+1}=\frac{1}{1+x_n}\leqslant\frac{2}{3}\)，仍在区间 \(\left[\frac{1}{2},\frac{2}{3}\right]\) 内. 令 \(d_n=\left|x_{n+1}-x_n\right|\)，则 \(d_1=\left|\frac{2}{3}-\frac{1}{2}\right|=\frac{1}{6}\). 当 \(n\geqslant2\) 时，
\[
d_n=\left|\frac{1}{1+x_n}-\frac{1}{1+x_{n-1}}\right|=\frac{d_{n-1}}{\left(1+x_n\right)\left(1+x_{n-1}\right)}.
\]
又因为 \(x_n=\frac{1}{1+x_{n-1}}\)，所以 \(\left(1+x_n\right)\left(1+x_{n-1}\right)=2+x_{n-1}\geqslant\frac{5}{2}\)，从而 \(d_n\leqslant\frac{2}{5}d_{n-1}\). 反复使用并结合 \(d_1=\frac{1}{6}\)，得 \(d_n\leqslant\frac{1}{6}\left(\frac{2}{5}\right)^{n-1}\)，即 \(\left|x_{n+1}-x_n\right|\leqslant\frac{1}{6}\left(\frac{2}{5}\right)^{n-1}\).
\end{enumerate}
\end{solution}
